\documentclass[11pt]{article}
\usepackage[margin=2.5cm]{geometry}
\usepackage{amsmath,amssymb}
\usepackage{booktabs}
\usepackage[T1]{fontenc}
\usepackage{lmodern}

\newcommand{\yes}{\text{yes}}
\newcommand{\no}{\text{no}}
\newcommand{\tab}{\text{tab}}
\newcommand{\code}[1]{\texttt{#1}}

\title{A model of the parser of makefile-doc.awk}
\date{}

\begin{document}
\maketitle

\section{Purpose}

The parser reads the lines of one or more Makefiles and attaches each doc (and section) to
its anchor: a target or a variable. This document models the parser as a state machine: a
state that evolves line by line, driven by an observation of each line. It is the
specification for the redesign of the parser.

\section{Assumptions}

\begin{enumerate}
  \item \textbf{Logical lines are given.} Before the model runs, every physical line that
        ends with an odd number of backslashes is joined with the next one, recipe lines
        included. The model sees the resulting logical lines $l_1, l_2, \dots$, over all
        files in order.
  \item \textbf{Files are read whole, one after another}, in the order given (for
        \code{make help}, the order of \code{\$(MAKEFILE\_LIST)}).
\end{enumerate}

\section{Notation}

\begin{center}
\begin{tabular}{@{}ll@{}}
  \toprule
  Symbol & Description \\
  \midrule
  $l_i$ & the $i$-th logical line \\
  $b_i$ & line $i$ is the first line of a file: \yes/\no{} (an input, not read from the line) \\
  $r_i$ & a rule is open from line $i$: \yes/\no \\
  $d_i$ & a doc is pending from line $i$: \yes/\no \\
  $s_i$ & a section is pending from line $i$: \yes/\no \\
  $c_i$ & the recipe prefix from line $i$ (a character) \\
  $x_i$ & the state, $(r_i, d_i, s_i, c_i)$ \\
  $K$ & the finite set of kinds of line \\
  $k_i$ & the kind of line $i$, $k_i \in K$ \\
  $a_i$ & line $i$ has an inline doc: \yes/\no \\
  \bottomrule
\end{tabular}
\end{center}

\noindent ``From line $i$'' means: starting at line $i$, with the effect of line $i$
included, until a later line changes it.

\section{Model}

\begin{align*}
  x_0 &= (\no,\ \no,\ \no,\ c_0)
  \qquad \text{($c_0$: the \code{RECIPEPREFIX} option, or a tab)} \\[6pt]
  r_{i-1} &\leftarrow \begin{cases} \no & \text{if } b_i \\ r_{i-1} & \text{otherwise} \end{cases}
  \qquad
  d_{i-1} \leftarrow \begin{cases} \no & \text{if } b_i \\ d_{i-1} & \text{otherwise} \end{cases} \\[6pt]
  (k_i,\ a_i,\ c_i) &= \mathrm{obs}(l_i,\ r_{i-1},\ c_{i-1}) \\[2pt]
  (r_i,\ d_i,\ s_i) &= \mathrm{step}(r_{i-1},\ d_{i-1},\ s_{i-1},\ k_i,\ a_i)
\end{align*}

\begin{itemize}
  \item \textbf{The reset at a file boundary} ($\leftarrow$: an update applied before
        line $i$ is processed) clears $r$ and $d$, and keeps $s$ (sections carry over to
        the next file) and $c$ (GNU Make does not reset \code{.RECIPEPREFIX} between files).
  \item \textbf{$\mathrm{obs}$} reads the line. It needs $r_{i-1}$ and $c_{i-1}$ for one
        decision: a line is a recipe line if a rule is open and the line starts with the
        recipe prefix. It returns $c_i$: the new prefix if $l_i$ assigns
        \code{.RECIPEPREFIX} (a tab if the value is empty), $c_{i-1}$ otherwise.
  \item \textbf{$\mathrm{step}$} updates $r$, $d$ and $s$ from the observation alone; it
        never reads the line.
\end{itemize}

\section{Output}

$\mathrm{step}$ produces the attachments: line $i$ attaches a doc when it is an anchor and
$d_{i-1}$ or $a_i$ holds (a pending doc takes precedence over an inline doc), with the
pending section if $s_{i-1}$ holds. The content (doc text, section text, anchor names,
tags) travels alongside the state as data: control depends only on whether something is
pending.

Postprocessing, outside the model: the warnings (redefined docs, docs pending at the end
of a file, a continued last line) and the names of double-colon rules (\code{t\textasciitilde 1},
\code{t\textasciitilde 2}).

\section{Decisions}

Where the model deliberately differs from the current parser:
\begin{itemize}
  \item $c$ is tracked: an assignment to \code{.RECIPEPREFIX} changes the prefix of the
        lines that follow (today, the prefix is the \code{RECIPEPREFIX} option, constant).
  \item Recipe continuations are joined like the other lines. This changes the
        representation only: where a recipe line ends is the same.
  \item A line that holds only the recipe prefix (possibly followed by blanks) while a rule
        is open is a recipe line, as in GNU Make (today, it is a blank line: pending docs
        survive it).
  \item $\mathrm{obs}$ follows GNU Make, which fixes issues \#181, \#182 and \#183.
\end{itemize}

\section{Limits}

Where the model cannot follow GNU Make:
\begin{itemize}
  \item An included file is read by GNU Make at the \code{include} line, in the middle of
        its parent; the model reads the files one after another.
  \item Some values of \code{.RECIPEPREFIX} cannot be known without evaluating the
        Makefile: an expansion (\code{.RECIPEPREFIX := \$(P)}), an assignment inside a
        conditional, an assignment in an included file. $\mathrm{obs}$ then keeps
        $c_{i-1}$ and reports a warning.
  \item Conditionals are not evaluated: the lines of all their branches are read in order.
\end{itemize}

\section{To be specified}

The set of kinds $K$ with the rule that recognizes each kind (the core of $\mathrm{obs}$),
and, for each kind, its effect on $r$, $d$ and $s$ (the core of $\mathrm{step}$).

\end{document}
